\documentclass[../../main.tex]{subfiles}
\begin{document}
\section{de Rham 上同调群}\label{section:ch17-the-de-rham-cohomology-groups}

在第~14~章中我们定义了闭形式与恰当形式：一个光滑微分形式 $\omega$ 若满足 $\mathrm{d}\omega=0$ 则称为闭的；若它可以写成 $\omega=\mathrm{d}\eta$ 则称为恰当的。由于 $\mathrm{d}\circ\mathrm{d}=0$，每个恰当形式都是闭的。在本章中，我们探讨反方向问题的含义：是否每个闭形式都是恰当的？一般答案是否定的：例如在例~11.48~中，我们看到 punctured 平面上有一个 $1$-形式是闭的但不是恰当的；恰当性的失效似乎是域中心处“洞”的一个后果。对于更高次的形式，哪些闭形式是恰当的这一问题，依赖于流形微妙的拓扑性质，这些性质与更高维“洞”的存在有关。把这种依赖定量化，就导出光滑流形的一组新的不变量，称为 de Rham 上同调群，这正是本章的主题。

知道哪些闭形式是恰当的有许多重要推论。例如，Stokes 定理表明：若 $\omega$ 是恰当的，则 $\omega$ 在任意不带边紧子流形上的积分为零。

设 $M$ 是一个带边或不带边的光滑流形，$p$ 是非负整数。由于 $\mathrm{d}:\Omega^p(M)\to\Omega^{p+1}(M)$ 是线性的，其核与像都是线性子空间。定义
\begin{align*}
Z^p(M)&=\operatorname{Ker}\bigl(\mathrm{d}:\Omega^p(M)\to\Omega^{p+1}(M)\bigr)=\{M\text{ 上的闭 }p\text{-形式}\},\\
B^p(M)&=\operatorname{Im}\bigl(\mathrm{d}:\Omega^{p-1}(M)\to\Omega^p(M)\bigr)=\{M\text{ 上的恰当 }p\text{-形式}\}.
\end{align*}
按约定，当 $p<0$ 或 $p>n=\dim M$ 时 $\Omega^p(M)$ 视为零向量空间。每个恰当形式都是闭的这一事实意味着 $B^p(M)\subseteq Z^p(M)$。于是可以合理地把 $M$ 的 $p$ 次 de Rham 上同调群（或第 $p$ 个 de Rham 群）定义为商向量空间
\begin{align}
H^p_{\mathrm{dR}}(M)=Z^p(M)/B^p(M).
\end{align}
（它是一个实向量空间，特别地在向量加法下是一个群。也许“de Rham 上同调空间”是一个更恰当的名称，但由于大多数其他上同调理论只产生群，传统上在此情形也沿用“群”这一名称，要记住这些“群”实际上是实向量空间。）显然当 $p<0$ 或 $p>\dim M$ 时 $H^p_{\mathrm{dR}}(M)=0$。对 $0\leqslant p\leqslant n$，定义表明 $H^p_{\mathrm{dR}}(M)=0$ 当且仅当 $M$ 上每个闭 $p$-形式都是恰当的。

\begin{example}\label{example:ch17-punctured-plane-not-exact}
$\mathbb{R}^2\setminus\{0\}$ 上存在一个闭但不恰当的 $1$-形式，这一事实意味着 $H^1_{\mathrm{dR}}(\mathbb{R}^2\setminus\{0\})\neq0$（见例~11.48）。另一方面，$1$-形式的 Poincaré 引理（定理~11.49）表明，对任意星形开子集 $U\subseteq\mathbb{R}^n$，$H^1_{\mathrm{dR}}(U)=0$。
\end{example}

首先要做的是证明 de Rham 群是微分同胚不变量。对 $M$ 上任意闭 $p$-形式 $\omega$，记 $[\omega]$ 为 $\omega$ 在 $H^p_{\mathrm{dR}}(M)$ 中的等价类，称为 $\omega$ 的\emph{上同调类}。若 $[\omega]=[\omega']$（即 $\omega$ 与 $\omega'$ 相差一个恰当形式），则称 $\omega$ 与 $\omega'$ 上同调。

\begin{proposition}[诱导的上同调映射]\label{proposition:ch17-induced-cohomology-maps}
对任意带边或不带边光滑流形之间的光滑映射 $F:M\to N$，拉回 $F^*:\Omega^p(N)\to\Omega^p(M)$ 把 $Z^p(N)$ 映入 $Z^p(M)$，把 $B^p(N)$ 映入 $B^p(M)$。因此它诱导出一个线性映射（仍记为 $F^*$）$F^*:H^p_{\mathrm{dR}}(N)\to H^p_{\mathrm{dR}}(M)$，称为\emph{诱导上同调映射}。它有下列性质：
\begin{enumerate}[(a)]
\item 若 $G:N\to P$ 是另一个光滑映射，则 $(G\circ F)^*=F^*\circ G^*$。
\item 若 $\operatorname{Id}$ 是 $M$ 的恒等映射，则 $\operatorname{Id}^*$ 是 $H^p_{\mathrm{dR}}(M)$ 的恒等映射。
\end{enumerate}
\end{proposition}

\begin{proof}
若 $\omega$ 闭，则 $\mathrm{d}(F^*\omega)=F^*(\mathrm{d}\omega)=0$，故 $F^*\omega$ 也闭。若 $\omega=\mathrm{d}\eta$ 恰当，则 $F^*\omega=F^*(\mathrm{d}\eta)=\mathrm{d}(F^*\eta)$，也恰当。因此 $F^*$ 把 $Z^p(N)$ 映入 $Z^p(M)$，把 $B^p(N)$ 映入 $B^p(M)$。诱导上同调映射定义为：对闭 $p$-形式 $\omega$，令 $F^*[\omega]=[F^*\omega]$。函子性质由拉回的相应函子性质得到。
\end{proof}

\begin{proposition}[不相交并的 de Rham 群]\label{proposition:ch17-disjoint-union}
设 $\{M_j\}$ 是任意一族带边或不带边光滑流形，$M=\coprod_j M_j$。则包含映射 $\iota_j:M_j\hookrightarrow M$ 诱导同构
\begin{align*}
H^p_{\mathrm{dR}}(M)\cong\prod_j H^p_{\mathrm{dR}}(M_j).
\end{align*}
\end{proposition}

\begin{proposition}[连通分支上的 $0$ 次群]\label{proposition:ch17-zero-degree}
设 $M$ 是一个光滑流形，其连通分支为 $\{M_j\}$。则 $H^0_{\mathrm{dR}}(M)$ 是 $M$ 上局部常值实值函数的向量空间，并且当 $M$ 连通时 $H^0_{\mathrm{dR}}(M)=\mathbb{R}$。
\end{proposition}
\end{document}
